% Stage 8 restored fragment — 2.7.3 Evanescent overlap and non-reconstructability (archive lines 9570-9849)
\subsection{Evanescent overlap and non-reconstructability}
Subsections~\ref{sec:ch271}--\ref{sec:ch272} fixed $\ker\mathcal{R}_\omega$ and canonical
non-radiating impressed currents Eq.~\eqref{eq:ch2.nr-current-def}. Evanescent sheets are a
distinct field-side deficit: reactive storage that never passes $\Pi_{\mathrm{rad}}$ even when
the source radiates. Subsection~\ref{sec:ch273} supplies the definitive account of evanescent
subspaces, their overlap with non-radiating source content, and reconstructability limits under
finite observation rank---developing Eq.
Eq.~\eqref{eq:ch2.evanescent-nonreconstruct} from Section~\ref{sec:ch27}
\cite{stratton1941}. Propagation projectors
Eqs.~\eqref{eq:ch2.spectral-projectors}--\eqref{eq:ch2.rad-null-on-evanescent} from
Subsection~\ref{sec:ch224}, non-propagating split
Eq.~\eqref{eq:ch2.nonpropagating-component-def} from Subsection~\ref{sec:ch254}, measurement
null-space extension measurement null-space extension (Subsection~\ref{sec:ch271}) from
Subsection~\ref{sec:ch271}, and rank cap Eq.~\eqref{eq:ch2.reconstruction-rank-cap} from
Subsection~\ref{sec:ch244} are presupposed; projector algebra and inverse-source templates are
not re-proved \cite{stratton1941}. Angular-spectrum consequences belong to Subsection
~\ref{sec:ch274}; plane-wave filters to Section~\ref{sec:ch29} \cite{harrington2001}.
Assumptions. Fix \(\omega>0\), admissible \(\mathcal{J}_\Omega\) Eq.~\eqref{eq:ch2.source-admissibility},
outgoing radiation condition Eqs.~\eqref{eq:ch2.outgoing-radiation-refinement}--\eqref{eq:ch2.outgoing-resolvent-unique},
field pairs in \(\mathcal{F}_{\mathrm{SF}}\subset\mathcal{F}_\Omega\), and spectral completion
supporting Eqs.~\eqref{eq:ch2.spectral-projectors}--\eqref{eq:ch2.spectral-projector-algebra}
\cite{stratton1941}. Observation map
\(\mathcal{M}_\omega=\mathcal{O}_{\mathrm{meas}}\circ\Lambda_{\mathrm{rad}}\) from
Eq.~\eqref{eq:ch2.reconstruction-observation-map} is finite rank with
\(\mathrm{rank}(\mathcal{O}_{\mathrm{meas}})=N\) Eq.~\eqref{eq:ch2.finite-rank-observation}
\cite{stratton1941}. regime qualification Eq.~\eqref{eq:ch2.regime-operator-screen} governs transport
reports on any recovered radiative component \cite{harrington2001}. On the spectral completion \(\widehat{\mathcal{F}}_\Omega\), the \emph{evanescent subspace} is
\begin{equation}
\label{eq:ch2.evanescent-subspace-def}
\mathcal{F}_{\mathrm{ev}}
=
\left\{
\mathbf{u}\in\widehat{\mathcal{F}}_\Omega:
\Pi_{\mathrm{ev}}\mathbf{u}=\mathbf{u}
\right\},
\end{equation}
with \(\Pi_{\mathrm{ev}}\) from Eq.~\eqref{eq:ch2.spectral-projectors}
\cite{stratton1941}. Equation~\eqref{eq:ch2.evanescent-subspace-def} is the
Subsection~\ref{sec:ch273} field-side label for modes with imaginary longitudinal
wavenumber Eq.~\eqref{eq:ch2.longitudinal-wavenumber-split} and decay rate
Eq.~\eqref{eq:ch2.evanescent-decay-rate} \cite{stratton1941}. The companion \emph{non-propagating} subspace combines evanescent and
below-cutoff content through Eq.~\eqref{eq:ch2.nonpropagating-component-def}
\cite{harrington2001}. Evanescent sheets carry reactive structure bound to apertures; they decay
exponentially with scan height and satisfy Eq.~\eqref{eq:ch2.rad-null-on-evanescent} without
re-derivation here \cite{stratton1941}.
For admissible \(\Jimp\in\mathcal{J}_\Omega\), the evanescent projector \(\Pi_{\mathrm{ev}}\)
selects field content whose lateral wavenumber exceeds the vacuum propagation circle
\(k_t>k_0\): such modes decay exponentially with distance from the source and carry reactive
storage rather than qualified export. Define the evanescent field content
\begin{equation}
\label{eq:ch2.evanescent-source-overlap}
\mathbf{u}_{\mathrm{ev}}
=
\Pi_{\mathrm{ev}}\,\mathcal{S}_\omega[\Jimp],
\qquad
\|\mathbf{u}_{\mathrm{ev}}\|_{\mathcal{E}}
=
\|\Pi_{\mathrm{ev}}\mathcal{S}_\omega[\Jimp]\|_{\mathcal{E}},
\end{equation}
with energy norm Eq.~\eqref{eq:ch2.field-energy-norm} \cite{stratton1941}.
Equation~\eqref{eq:ch2.evanescent-source-overlap} measures reactive storage excited by
\(\Jimp\) that never passes through \(\Pi_{\mathrm{rad}}\) because
\begin{equation}
\label{eq:ch2.evanescent-nonexport-law}
\Pi_{\mathrm{rad}}\,\mathbf{u}_{\mathrm{ev}}=\mathbf{0},
\qquad
\mathcal{R}_\omega[\Jimp]\ \text{depends only on}\ \Pi_{\mathrm{prop}}\mathcal{S}_\omega[\Jimp]
\ \text{after outgoing selection},
\end{equation}
by Eqs.~\eqref{eq:ch2.rad-null-on-evanescent} and~\eqref{eq:ch2.rad-projection-composition}
\cite{stratton1941}. Equation~\eqref{eq:ch2.evanescent-nonexport-law} upgrades the second
line of Eq.~\eqref{eq:ch2.evanescent-nonreconstruct}: evanescent overlap is
field-side non-export, distinct from source-side \(\Jimp\in\ker\mathcal{R}_\omega\)
Eq.~\eqref{eq:ch2.radiation-nullspace-def} \cite{harrington2001}. A source may satisfy
\(\|\mathbf{u}_{\mathrm{ev}}\|\gg\|\mathcal{R}_\omega[\Jimp]\|\) while still radiating
\cite{stratton1941}.
Eq.~\eqref{eq:ch2.evanescent-nonreconstruct} separates export annihilation from
finite-rank readout loss. Source nullity \(\ker\mathcal{R}_\omega\) is strict export failure;
measurement nullity \(\ker\mathcal{M}_\omega\) additionally loses directions that radiate but
the aperture cannot resolve \cite{stratton1941}:
\begin{equation}
\label{eq:ch2.radiation-vs-measurement-nullity-split}
\underbrace{\ker\mathcal{R}_\omega}_{\text{no export}}
\ \subseteq\
\underbrace{\ker\mathcal{M}_\omega}_{\text{no finite readout}},
\qquad
\ker\mathcal{M}_\omega
\supsetneq
\ker\mathcal{R}_\omega
\ \text{when}\ 
\ker\mathcal{O}_{\mathrm{meas}}\big|_{\operatorname{ran}\mathcal{R}_\omega}\neq\{\mathbf{0}\},
\end{equation}
with strict inclusion typical for finite apertures
Eq.~\eqref{eq:ch2.observation-null-space}
\begin{equation}
\label{eq:ch2.measurement-nullspace-composition}
\ker\mathcal{M}_\omega
=
\ker\mathcal{R}_\omega
\ \cup\
\mathcal{R}_\omega^{-1}\!\big(\ker\mathcal{O}_{\mathrm{meas}}\big|_{\operatorname{ran}\mathcal{R}_\omega}\big),
\end{equation} \cite{stratton1941}. Equation
\eqref{eq:ch2.radiation-vs-measurement-nullity-split} expresses measurement null-space extension (Subsection~\ref{sec:ch271}) and
Eq.~\eqref{eq:ch2.measurement-nullspace-composition} without re-proof: evanescent storage
contributes to the first factor only when it annihilates export; finite rank contributes to
the second factor even when export is nonzero \cite{stratton1941}. Tikhonov reconstruction
Eq.~\eqref{eq:ch2.tikhonov-source-reconstruction} stabilizes ill-conditioned inversions on
\(\operatorname{ran}\mathcal{M}_\omega^{\dagger}\) but does not shrink
\(\ker\mathcal{M}_\omega\) \cite{harrington2001}.
Define the \emph{non-reconstructable set} at observation rank \(N\):
\begin{equation}
\label{eq:ch2.nonreconstruct-observation-kernel}
\mathcal{K}_{\mathrm{nrec}}(N)
=
\left\{
\Jimp\in\mathcal{J}_\Omega:
\mathcal{M}_\omega[\Jimp]=\mathbf{0}
\right\}
=
\ker\mathcal{M}_\omega,
\end{equation}
equivalently the kernel of Eq.~\eqref{eq:ch2.inverse-source-equation} for generic
\(\mathbf{y}\) \cite{stratton1941}. Equation
\eqref{eq:ch2.nonreconstruct-observation-kernel} is the Subsection~\ref{sec:ch273} constructive account of
reconstructability label: sources in \(\mathcal{K}_{\mathrm{nrec}}(N)\) are invisible to the
declared measurement bank regardless of near-zone \(|\E|\) \cite{stratton1941}. Rank bound
\begin{equation}
\label{eq:ch2.reconstructability-rank-law}
\dim\operatorname{ran}\mathcal{M}_\omega
\le
N
<
\infty
\ \Longrightarrow\
\dim\mathcal{K}_{\mathrm{nrec}}(N)=\infty
\ \text{on typical open domains},
\end{equation}
follows from Eq.~\eqref{eq:ch2.reconstruction-rank-cap} and infinite-dimensional
\(\mathcal{J}_\Omega\) \cite{stratton1941}. Geometric sector constraints
Eq.~\eqref{eq:ch2.accessible-modal-sector} and composition law
Eq.~\eqref{eq:ch2.constraint-composition-law} further restrict recoverable coordinates
independently of Eq.~\eqref{eq:ch2.nonreconstruct-observation-kernel}
\cite{stratton1941}.
\begin{figure}[t]
\centering
\ChOneFigBlock{%
\begin{tikzpicture}[ch1fig, node distance=5mm and 7mm]
 \node[ch1box] (j) {$\mathcal{J}_\Omega$};
 \node[ch1box, right=7mm of j] (s) {$\mathcal{S}_\omega$};
 \node[ch1box, right=7mm of s] (ev) {$\Pi_{\mathrm{ev}}$};
 \node[ch1box, right=7mm of ev] (r) {$\mathcal{R}_\omega$};
 \node[ch1box, right=7mm of r] (m) {$\mathcal{M}_\omega$};
 \node[ch1box, below=5mm of ev] (uev) {$\mathbf{u}_{\mathrm{ev}}$};
 \node[ch1box, below=5mm of m] (kn) {$\mathcal{K}_{\mathrm{nrec}}$};
 \draw[ch1flow] (j) -- (s) -- (ev) -- (r) -- (m);
 \draw[ch1flow] (ev.south) -- (ev.south |- uev.north);
 \draw[ch1flow] (ev.south |- uev.north) -- (uev.north);
 \draw[ch1flow] (m.south) -- (m.south |- kn.north);
 \draw[ch1flow] (m.south |- kn.north) -- (kn.north);
\end{tikzpicture}%
}
\caption{Subsection~\ref{sec:ch273} measurement sequence: evanescent overlap
Eq.~\eqref{eq:ch2.evanescent-source-overlap} never exports through
Eq.~\eqref{eq:ch2.evanescent-nonexport-law}; finite-rank observation maps
\(\mathcal{J}_\Omega\) to \(\mathcal{K}_{\mathrm{nrec}}(N)\)
Eq.~\eqref{eq:ch2.nonreconstruct-observation-kernel}.}
\label{fig:ch2.evanescent-nonreconstruct-measurement chain}
\end{figure}
The measurement sequence orders radiation selection, evanescent overlap, and finite-rank observation. $\Pi_{\mathrm{ev}}$ content may dominate $|\E|$ on inner shells yet satisfies Eq.~\eqref{eq:ch2.evanescent-nonexport-law} on qualified flux meters. $\mathcal{O}_{\mathrm{meas}}$ with rank $N$ collapses recoverable modes into $\mathcal{K}_{\mathrm{nrec}}(N)$: export can exist while the aperture cannot resolve it. Dashed branches mark observation deficits distinct from source nullity $\ker\mathcal{R}_\omega$ \cite{stratton1941}. \phantomsection\label{thm:ch2.evanescent-nonreconstructability}
Let \(\Jimp\in\mathcal{J}_\Omega\) be admissible with evanescent content
Eq.~\eqref{eq:ch2.evanescent-source-overlap} and observation map
\(\mathcal{M}_\omega\) of rank \(N\). Then:
\begin{enumerate}
\item[(i)] \(\Pi_{\mathrm{rad}}\mathbf{u}_{\mathrm{ev}}=\mathbf{0}\) and
\(\mathbf{u}_{\mathrm{ev}}\) carries no qualified far-zone export
Eq.~\eqref{eq:ch2.evanescent-nonexport-law}.
\item[(ii)] \(\ker\mathcal{R}_\omega\subseteq\ker\mathcal{M}_\omega=\mathcal{K}_{\mathrm{nrec}}(N)\).
\item[(iii)] Near-field data on a finite enclosing surface reduce but do not in general
eliminate \(\mathcal{K}_{\mathrm{nrec}}(N)\): if \(\mathcal{M}_{\mathrm{near}}\) samples
\(\mathcal{S}_\omega[\Jimp]\) on \(\Gamma_{\mathrm{enc}}\), then
\begin{equation}
\label{eq:ch2.near-field-nullity-reduction}
\ker\mathcal{M}_{\mathrm{near}}
\subsetneq
\ker\mathcal{M}_\omega
\ \text{when}\ \Gamma_{\mathrm{enc}}\ \text{resolves more modes than far-zone aperture},
\end{equation}
while \(\ker\mathcal{R}_\omega\) is unchanged \cite{stratton1941}.
\item[(iv)] Reconstructability budget:
\begin{equation}
\label{eq:ch2.reconstructability-budget}
\mathcal{N}_{\mathrm{reportable}}
\subseteq
\mathcal{N}_{\mathrm{acc}}
\cap
\{n:\ \text{resolvable by}\ \mathcal{M}_\omega\}
\cap
\{n:\ \lambda_n\neq 0\},
\end{equation}
with \(\mathcal{N}_{\mathrm{acc}}\) from Eq.~\eqref{eq:ch2.accessible-modal-sector} and modal
indices from Subsection~\ref{sec:ch261} \cite{stratton1941}.
\end{enumerate}
Part (i) is Eq.~\eqref{eq:ch2.rad-null-on-evanescent} \cite{stratton1941}. Part (ii) follows
from measurement null-space extension (Subsection~\ref{sec:ch271}) and definition
Eq.~\eqref{eq:ch2.nonreconstruct-observation-kernel} \cite{harrington2001}. Part (iii) uses
near- versus far-field sampling logic from Subsection~\ref{sec:ch244}: \(\mathcal{S}_\omega\) is
single-valued on outgoing fields while \(\mathcal{R}_\omega\) remains non-injective
\cite{stratton1941}. Part (iv) composes
Eq.~\eqref{eq:ch2.constraint-composition-law} with rank cap
Eq.~\eqref{eq:ch2.reconstruction-rank-cap} \cite{stratton1941}.
\phantomsection\label{prop:ch2.evanescent-never-exports}
For all admissible \(\Jimp\),
\begin{equation}
\label{eq:ch2.evanescent-radiation-annihilation}
\mathcal{R}_\omega[\Jimp]
=
\Pi_{\mathrm{rad}}\,\mathcal{S}_\omega[\Jimp]
=
\Pi_{\mathrm{rad}}\,\Pi_{\mathrm{prop}}\,\mathcal{S}_\omega[\Jimp],
\end{equation}
and \(\Pi_{\mathrm{rad}}\Pi_{\mathrm{ev}}=\mathbf{0}\) on \(\mathcal{F}_\Omega\)
\cite{stratton1941}. Equation~\eqref{eq:ch2.evanescent-radiation-annihilation}
implies \(\mathbf{u}_{\mathrm{ev}}\) from Eq.~\eqref{eq:ch2.evanescent-source-overlap} does
not alter \(\mathcal{R}_\omega[\Jimp]\); adding evanescent excitation to an otherwise
radiating source changes \(\mathbf{u}_{\mathrm{nr}}\) and local power
Eq.~\eqref{eq:ch2.component-power-budget} without changing far patterns when export is
unchanged \cite{stratton1941}. regime qualification Eq.~\eqref{eq:ch2.regime-component-criterion} still
forbids reporting \(\mathbf{u}_{\mathrm{ev}}\) as transport \cite{harrington2001}. Substitute Eqs.~\eqref{eq:ch2.rad-null-on-evanescent} and~\eqref{eq:ch2.rad-projection-composition}
into \(\mathcal{R}_\omega=\Pi_{\mathrm{rad}}\circ\mathcal{S}_\omega\)
\cite{stratton1941}. Near-zone energy follows from
Eq.~\eqref{eq:ch2.nonpropagating-component-def} \cite{harrington2001}.
\phantomsection\label{prop:ch2.near-field-nullity-reduction}
Let \(\Gamma_{\mathrm{enc}}\) be a closed surface enclosing \(\Omega_s\) and let
\(\mathcal{M}_{\mathrm{near}}\) sample tangential \(\E\), \(\H\) on \(\Gamma_{\mathrm{enc}}\).
Then \(\ker\mathcal{M}_{\mathrm{near}}\subseteq\ker\mathcal{M}_\omega\) when far-zone
observations are a subset of near-field functionals, but
\begin{equation}
\label{eq:ch2.near-field-kernel-inclusion}
\ker\mathcal{R}_\omega
\subseteq
\ker\mathcal{M}_{\mathrm{near}}
\subseteq
\ker\mathcal{M}_\omega
\end{equation}
with strict inclusions possible at each step \cite{stratton1941}. Equation
\eqref{eq:ch2.near-field-kernel-inclusion} expresses Subsection~\ref{sec:ch244} near- versus
far-field data without re-opening uniqueness exterior radiative uniqueness (Subsection~\ref{sec:ch243}) \cite{stratton1941}. Outgoing radiation condition remains
mandatory: identical near traces on a finite surface do not imply identical far zones without
exterior selection \cite{harrington2001}.
If \(\mathcal{M}_\omega[\Jimp]=\mathbf{0}\), then far functionals vanish; near samples may
still vanish on \(\Gamma_{\mathrm{enc}}\) for non-radiating additions in
\(\ker\mathcal{R}_\omega\) \cite{stratton1941}. Additional near resolution shrinks the
observation null space on \(\operatorname{ran}\mathcal{R}_\omega\) but leaves \(\ker\mathcal{R}_\omega\) unchanged Eq.~\eqref{eq:ch2.radiation-nullspace-def}
\cite{harrington2001}.
When Eq.~\eqref{eq:ch2.cutoff-condition} fails, modal content lies in
\(\Pi_{\mathrm{non}}\mathcal{S}_\omega[\Jimp]\) or \(\Pi_{\mathrm{ev}}\mathcal{S}_\omega[\Jimp]\)
\cite{stratton1941}. Such storage satisfies
Eq.~\eqref{eq:ch2.evanescent-nonexport-law} and may dominate
Eq.~\eqref{eq:ch2.component-power-budget} on sub-wavelength shells while
\(\mathcal{P}_{\mathrm{rad}}=0\) \cite{stratton1941}. Applying far-zone pattern formulas to
below-cutoff modes violates Eq.~\eqref{eq:ch2.regime-operator-screen} as in the TE\(_{10}\)
illustration referenced from Section~\ref{sec:ch01.1.2} \cite{harrington2001}. Subsection
~\ref{sec:ch273} records operator consequences; guide nomenclature is not re-tabulated
\cite{stratton1941}. Non-radiating configurations Eq.~\eqref{eq:ch2.nr-current-def} may carry large
\(\|\mathbf{u}_{\mathrm{ev}}\|\) when \(\mathcal{S}_\omega[\Jimp^{\mathrm{(NR)}}]\) is
evanescent-dominated \cite{stratton1941}. Conversely, large
\(\|\mathbf{u}_{\mathrm{ev}}\|\) does not imply \(\Jimp\in\ker\mathcal{R}_\omega\): a radiating
source may store substantial evanescent energy near \(\Omega_s\) while exporting through
\(\Pi_{\mathrm{prop}}\) \cite{stratton1941}. Subsections~\ref{sec:ch272} and~\ref{sec:ch273}
address orthogonal deficits---source nullity versus reactive storage versus finite-rank readout
\cite{harrington2001}.
Subsection~\ref{sec:ch273} assumes linear Maxwell theory at fixed
\(\omega\), outgoing radiation condition, and declared finite-rank \(\mathcal{O}_{\mathrm{meas}}\)
\cite{stratton1941}. Loss modifies evanescent decay
Eq.~\eqref{eq:ch2.evanescent-decay-rate} without promoting \(\Pi_{\mathrm{ev}}\) content to
\(\Pi_{\mathrm{rad}}\) export \cite{stratton1941}. Dispersive or broadband superposition
requires Subsection~\ref{sec:ch223}; harmonic evanescent labels here are fixed-frequency slices
\cite{harrington2001}. Subsection~\ref{sec:ch274} translates non-reconstructability to
continuous angular spectra Eq.~\eqref{eq:ch2.exterior-continuous-spectral-integral}
\cite{stratton1941}.
